\scanpage{6}
% archive foliation: 5
Pour le composé dans l'autre sens, on trouve que $G$ opère sur l'ens $I_G$
de ses \unsure{trois} ss-groupes de Sylow par autom. intérieurs, d'où
\[
 G\xrightarrow{\sim}\unsure{\mathcal{G}_{I_G}|\mathfrak{S}_{I_G}}
\]
qui est un iso.

Éléments de $\unsure{\mathcal{G}_I|\mathfrak{S}_I}$.

a) 1

b) Les \inserted{trois} éléments \deleted{non nuls de}\inserted{de} $V^*$,
qui forment une orbite $J$ sous $\unsure{\mathcal{G}_I|\mathfrak{S}_I}$, et correspondent
aux élém. de $\mathcal{P}_{2,2}(I)$.
\margin{``translations'' (non nulles) [Sym. p.r. aux trois axes du cube]}

c) Les \deleted{\unsure{trois|deux}} 8 éléments d'ordre 3, se répartissent en
\deleted{quatre}\inserted{deux} \deleted{torseurs} paquets de 4 qui sont des torseurs
(à gauche ou à droite au choix) sous $V$, suivant les deux éléments
$\ne1$ de $\unsure{\mathcal{G}_J|\mathfrak{S}_J}^+$.
L'ens de ces éléments s'identifie ainsi à
\[
 I\times\underline{\omega},\qquad
 \text{où }\underline{\omega}=\omega(J)\;(\simeq\omega(I))
\]
par l'application qui associe à tout
$\rho\in\unsure{\mathcal{G}_I|\mathfrak{S}_I}$ d'ordre 3 son unique pt fixe dans $I$
(ou encore, le sous-gr. de Sylow qu'il engendre), \uline{et} l'orientation
qu'il détermine \deleted{\unsure{ainsi}} ou encore son image dans
$\unsure{\mathcal{G}_J|\mathfrak{S}_J}^+$.
\margin{``rotations'' (ayant un pt fixe et un seul, lequel transforme $I$ en
vectoriel isom. à $V$) [``rotations-sommets'' du cube]}

d) Les \deleted{\unsure{8}}\inserted{6} éléments
\deleted{\unsure{de}}\inserted{d'ordre 2} de $\unsure{\mathcal{G}_I|\mathfrak{S}_I}$
qui s'envoient dans $\unsure{\mathcal{G}_J|\mathfrak{S}_J}^-$, en correspondance 1-1
avec l'ens $\mathcal{P}_2(I)$ --- ce sont les \uline{transpositions} dans $I$
qui se groupent en 3 \deleted{\unsure{paires}} paires, pour la transposition
$\tau_A$, $\tau_{A'}$ \deleted{\unsure{est}} $(A'=I\setminus A)$ est l'unique
\inserted{autre} élément d'ordre 2 tel que $\tau_A\tau_{A'}\in V^*$.
\margin{``réflexions'' (ayant du moins une \deleted{\unsure{paire}} droite de pts fixes)
[``rotations-arêtes'' du cube]}

e) Les 6 éléments d'ordre 4 de $\unsure{\mathcal{G}_I|\mathfrak{S}_I}$, correspondent aux
6 ordres circulaires sur $I$. Pour
\margin{[``rotations-faces'' du cube]}
